\documentclass[11pt, a4paper]{article}
\input{preamble.tex}
\switchlang{en}
\begin{document}

\title{1981 -- Xerox Daybreak Mouse}
\date{}
\maketitle
\selectlanguage{english}

In March 1973, Xerox Corporation announced the Xerox Alto computer, considered the first workstation (or personal computer), as well as the first computer with a graphical user interface \cite{wiki}. The workstation also came with the first mass-produced mouse in the history of computers, the Xerox Alto Mouse (fig. \ref{fig:XeroxDaybreakPic}). The original mouse based on two wheels was in several years upgraded and started to use a ball. But still, those mice weren't too reliable according to the quotes given in \cite{mouses}: they picked up dirt, quickly became gummed, and no longer controlled the cursor. Alto user had to unplug the mouse when that happened, put it into the ``Dead Mice''  box, and grab a cleaned mouse from the ``Clean Mice'' box. And the price of such a mouse was more than \$400.

The work for a completely optical mouse has been started near 1980 to solve two problems: to get higher reliability (because of no moving parts except the button switches), and much cheaper price (due to the one-chip design). The result of this effort was the last generation of the Alto mice -- Xerox Daybreak mouse, brought into life in 1981 \cite{vlsi81}. The mouse proved to be very successful compared to its mechanical predecessor, so its design was soon adopted for Xerox Star computers and later even for some of the company's copiers \cite{mouses}.

\begin{figure}[h]
    \centering
    \includegraphics[scale=0.7]{1985_xerox_daybreak_mouse/pic_30.jpg}
    \caption{Xerox Daybreak Mouse}
    \label{fig:XeroxDaybreakPic}
\end{figure}

\section*{Exterior}

The Xerox Daybreak mouse is housed in the same case as its mechanical version \cite{vlsi82}, and has the same color scheme. The case, made of (originally) cream-colored plastic, glossy in this mouse revision, is an almost perfect parallelepiped: it widens slightly toward the bottom and has convex sides to reduce the association with a ``box''. On the upper side of the case are three large, oblong, rounded buttons that meet at the edges, visually forming one solid block (Fig. \ref{XeroxDaybreakTopAndBottom}). Xerox software documentation called these buttons ``red'', ``yellow'', and ``blue''. However, in all cases (except for the earliest mouse modification with wheels and transversely oriented buttons \cite{vlsi81}), all three buttons were made of black (later, dark gray) plastic. It is likely that the myth that using one- and two-button mice was easier compared to a three-button one comes out of the confusion this unfortunate color coding caused among users.

\begin{figure}[h]
    \centering
    \includegraphics[scale=0.5]{1985_xerox_daybreak_mouse/top_30.jpg}
    \includegraphics[scale=0.5]{1985_xerox_daybreak_mouse/bottom_30.jpg}
    \caption{Xerox Daybreak Mouse, top and bottom views}
    \label{XeroxDaybreakTopAndBottom}
\end{figure}

A $4 \times 4$ size scanning elements matrix is responsible for reading color inhomogeneities at move, so the mouse required either a pad with a specially designed pattern or any surface with similar alternating light and dark spots, such as denim. The Xerox Daybreak mouse pad was a paper, sold in packs of 25 sheets \cite{pad}. The pattern was an array of light hexagons on a dark field and was easily reproduced by photocopying. A reconstructed pad can be seen in fig. \ref{fig:XeroxDaybreakPad}.

\begin{figure}[h]
    \centering
    \includegraphics[scale=0.3]{1985_xerox_daybreak_mouse/pad_30.jpg}
    \caption{Xerox Daybreak Mouse on a reconstructed pad}
    \label{fig:XeroxDaybreakPad}
\end{figure}

\section*{Size}

The mouse could have been smaller and lower in size if it had not used the standard Alto mouse body. As it is, it has the typical dimensions for mechanical mice of the first half of the 80s (figure \ref{fig:XeroxDaybreakSize}).

\begin{figure}[h]
    \centering
    \includegraphics[scale=0.4]{1985_xerox_daybreak_mouse/size_15.jpg}
    \caption{Xerox Daybreak Mouse on a graduated pad with a grid step of 1~cm}
    \label{fig:XeroxDaybreakSize}
\end{figure}

\section*{Ergonomics}

In terms of ergonomics and appearance, the Alto Mouse is minimalist. The user experience certainly suffers from the severe ``rectangularity'' of the case. Although this is partly compensated for by the convex oblong buttons, located within the reach of the fingers, the case doesn't provide much palm support (figure \ref{fig:XeroxDaybreakHand}).

\begin{figure}[h]
    \centering
    \includegraphics[scale=0.4]{1985_xerox_daybreak_mouse/hand_30.jpg}
    \caption{Xerox Daybreak Mouse with a human hand model}
    \label{fig:XeroxDaybreakHand}
\end{figure}

\section*{Connectivity}

The mouse has connector typical for Xerox Alto (DA-15 in this version), and the same quadrature interface as other Xerox mice.

\section*{Internal design}

Mouse internals are shown in figure \ref{fig:XeroxDaybreakInside}, where you can see the plastic cover which protects the scanning matrix from possible exposure to light (in later versions of mice for the Xerox Star, this cover was removed), as well as the signal processing chip. The printed circuit board is raised above the base, and angle-directed LEDs placed under are illuminating the mouse pad.

\begin{figure}[h]
    \centering
    \includegraphics[scale=0.8]{1985_xerox_daybreak_mouse/inside_30.jpg}
    \caption{Xerox Daybreak Mouse disassembled}
    \label{fig:XeroxDaybreakInside}
\end{figure}

\begin{thebibliography}{9}

\bibitem {vlsi81} R.\,F. Lyon. The Optical Mouse, and an Architectural Methodology for
Smart Digital Sensors // VLSI DESIGN, August 1981. - p. 20--30. \url{https://www.dicklyon.com/tech/OMouse/OpticalMouse-Lyon.pdf}
\bibitem {vlsi82} R.\,F. Lyon, M.\,P. Haeberli. Designing and Testing The Optical Mouse // VLSI DESIGN, January/February, 1982. - p. 20--30. \url{https://www.dicklyon.com/tech/OMouse/DesigningTestingOMouse.pdf}
\bibitem{pad} R.\,F. Lyon The Optical Mouse: Early Biomimetic Embedded Vision / Advances in Embedded Computer Vision, Nov 2014, pp.3-22 \url{https://static.googleusercontent.com/media/research.google.com/ru//pubs/archive/43260.pdf}
\bibitem{mouses} Xerox Mice. oldmouse.com \url{https://web.archive.org/web/20210418000634/http://oldmouse.com/mouse/xerox/alto.shtml}

\bibitem {star} Xerox Star: Wikipedia \url{https://en.wikipedia.org/wiki/Xerox_Star}

\bibitem{daybreak} Xerox Daybreak: Wikipedia \url{https://en.wikipedia.org/wiki/Xerox_Daybreak}

It was sold as the Xerox 6085 PCS (Professional Computer System) or ViewPoint 6085 PCS when sold as an office workstation running the ViewPoint system. ViewPoint is based on the Star software originally developed for the Xerox Star. 
The Daybreak was also sold as a Xerox 1186 workstation when configured as a Lisp machine.


\bibitem {photo} Alves P. Xerox 6085 workstation from 1985. Retro Microcomputers, Workstations, Servers and Consoles - Welcome! November 27, 2018. \url{https://www.facebook.com/groups/retrocomputers/posts/2333452826684585/}

photo of a computer, mouse, keyboard and printed mousepad on paper

\bibitem{starmouse} Degnan B. Xerox 8010 Star. - vintagecomputer.net 11/13/2011\url{https://vintagecomputer.net/browse_thread.cfm?id=311}

photo of original star mouse body

\bibitem{hist} GUI Wonderland #2: A Galaxy of Possibilities – Xerox Star & Daybreak. - The Authopire Compendium, April 16th 2024 \url{https://blisscast.wordpress.com/2024/04/16/xerox-star-daybreak-gui-wonderland-2/}

in 1985 Xerox released the very first version of the so-called ViewPoint software, dubbed 1.0, along with a new machine, the 6085 workstation, also known as Daybreak (we’ll have a look at all the several prototypes and hardware revisions of the Star in their own section), intended as a successor to the Star

the center button was removed to make the mouse less confusing, as the users of the time weren’t sure about the exact purpose of all those button

Eventually, in 1985, they also released the Daybreak system (the 6085 workstation), the successor to the Star that we mentioned before, which runs ViewPoint (more on that later on); its basic configuration had a 1.1 MB of RAM (upgradable to 3.7 MB), a 10 MB hard disk (upgradable to 80 MB), and a 15 or 19-inch monochrome display, and cost US$4,995 ($17,162.80 in 2024). Sadly, it never got a successor, marking the end of the workstation line started by the Alto.


\bibitem{office} Linden T., Harslem E. (editors). Office Systems Technology. A Look into the World of the Xerox 8000 Series Products: Workstations, Services, Ethernet, and Software Development, Xerox, November 1982. \ur{https://www.bitsavers.org/pdf/xerox/sdd/OSD-R8203_Xerox_Office_Systems_Technology_Nov82.pdf}

Both Star and the Alto use a pointing device called the mouse (Figure 2). First developed at SRI,7 Xerox's version has a ball on the bottom that turns as the mouse slides over a flat surface such as a table. Electronics sense the ball rotation and guide a cursor on the screen in corresponding motions.


There are two buttons on the mouse, called SELECT and ADJUST, used to make and adjust a selection 

\end{thebibliography}
\end{document}
